\documentclass[conference]{IEEEtran}
\IEEEoverridecommandlockouts
\usepackage[english]{babel}
\usepackage{cite}
\ifCLASSINFOpdf
 \usepackage[pdftex]{graphicx}
 \graphicspath{{./images/}{./svg-inkscape/}}
 \DeclareGraphicsExtensions{.pdf,.jpeg,.png}
\else
\fi
\usepackage{amsmath}
\usepackage{amssymb}
\newcommand{\argmax}{\operatorname*{argmax}}
\usepackage{array}
\usepackage{url}
\usepackage{hyperref}
\hyphenation{op-tical net-works semi-conduc-tor}
\usepackage{xcolor}
\usepackage{subcaption}
\usepackage{multirow}
\usepackage{placeins}

\title{Does Functional Similarity Induced by Distillation Preserve Internal Representations and Explanatory Mechanisms?\thanks{This research was funded by the Agencia Nacional de Investigaci\'on e Innovaci\'on (ANII), Uruguay, grant number POS\_NAC\_2025\_4\_189264.}}

\author{\IEEEauthorblockN{1\textsuperscript{st} Martin Luca Perciante}
\IEEEauthorblockA{\textit{dept. name of organization (of Aff.)} \\
\textit{Universidad Católica del Uruguay}\\
Montevideo, Uruguay \\
martinluca.perciante@ucu.edu.uy}
\and
\IEEEauthorblockN{2\textsuperscript{nd} Given Name Surname}
\IEEEauthorblockA{\textit{dept. name of organization (of Aff.)} \\
\textit{Universidad Católica del Uruguay}\\
Montevideo, Uruguay \\
email address or ORCID}
\and
\IEEEauthorblockN{3\textsuperscript{rd} Given Name Surname}
\IEEEauthorblockA{\textit{dept. name of organization (of Aff.)} \\
\textit{Universidad Católica del Uruguay}\\
Montevideo, Uruguay \\
email address or ORCID}
\and
\IEEEauthorblockN{4\textsuperscript{rd} Given Name Surname}
\IEEEauthorblockA{\textit{dept. name of organization (of Aff.)} \\
\textit{Universidad Católica del Uruguay}\\
Montevideo, Uruguay \\
email address or ORCID}
}

\begin{document}

\maketitle

\begin{abstract}
Knowledge distillation is often evaluated in terms of functional similarity between teacher and student models, measured through predictive performance. However, it remains unclear whether this functional similarity also implies the preservation of internal representations and explanatory mechanisms without explicitly imposing them. This work investigates this question in autoregressive Transformers. A \textbf{TinyLlama} \cite{zhang2024tinyllama} teacher is distilled into a family of 11 attention layer students with different training configurations. Internal representations are compared through layer-wise attention-map metrics, per-head similarity, while explanatory mechanisms are evaluated through \textbf{LIME}, \textbf{SHAP}, and \textbf{Integrated Gradients}, moving beyond qualitative examples to aggregate agreement and deletion-based faithfulness metrics. The study is conducted on binary \textbf{IMDb} sentiment classification and on four-class \textbf{AG News} topic classification. Results show that distillation does not transfer the teacher's explanatory mechanisms on either dataset. Representational similarity, in turn, is only partially and inconsistently attributable to distillation: attention-level and per-head metrics show that a non-distilled baseline is as close to the teacher as the distilled students. These findings suggest that the observed similarity to the teacher is driven more by the shared data and task than by distillation itself, indicating that functional similarity does not, by itself, guarantee the preservation of a teacher's explanatory mechanisms or internal representations.
\end{abstract}

\begin{IEEEkeywords}
Distillation, Attention, XAI, Explainability
\end{IEEEkeywords}

\section{Introduction}

Explainable Artificial Intelligence (XAI) seeks to interpret and explain the behavior of complex models, essential in domains where deployed systems must justify autonomous decisions and foster user trust. This task is particularly difficult for Large Language Models (LLMs), where probabilistic text generation replaces classification with a single class label output. Transformer-based models are the architecture of choice for this kind of modeling in current Natural Language Processing (NLP), driving recent progress largely by growing in parameters and complexity; the resulting computational cost has motivated compression techniques such as knowledge distillation (KD), in which a smaller model (student) is trained to approximate the output distribution of a larger model (teacher) \cite{hinton2015distilling}, alongside dataset training.

Standard KD is centered on functional similarity: the same input is made to produce a similar output, regardless of the internal representation that produce it. Prior work has shown that different attention configurations can induce similar functional behavior \cite{jain2019attention}, raising the question of whether KD enforces a preference for the teacher's internal representations, or whether it can converge to different attention maps while reproducing the teacher's output distribution. Bringing XAI into this picture, different models can provide different explanations for the same example while achieving similar performance over a dataset.

This work investigates whether the functional similarity induced by logit-based KD extends to similarity in internal representations and explanatory mechanisms. We adopt a generative classification setting, in which target classes correspond to single vocabulary tokens, allowing internal representations to be compared directly between teacher and student. We conducted the experiments on binary \textbf{IMDb} sentiment classification and in \textbf{AG News} four-class topic classification.

This work contributions are:
\begin{itemize}
    \item We show that logit-based KD, even when it reproduces the teacher's output distribution with increasing fidelity, does not correspondingly transfer its explanatory mechanisms.
    \item We show that representational similarity to the teacher (attention and per-head geometry) is not a reliable signature of distillation itself.
    \item We show this dissociation holds across both a binary (\textbf{IMDb}) and a four-class (\textbf{AG News}) classification task, indicating it is not an artifact of a single dataset or label space.
\end{itemize}

\section{Related Work}

Knowledge distillation (KD) transfers information from a high-capacity teacher to a smaller student by adding a Kullback--Leibler ($KL$) divergence term to the loss function, often softened by a temperature $T$ \cite{hinton2015distilling}. KD is widely used in deep networks and large language models, but its success is evaluated almost exclusively via functional metrics such as accuracy \cite{gou2021knowledge}. Some works explicitly bias the loss towards representational alignment, such as Attention Distillation \cite{zhou2025attention}, which directly supervises student attention maps against the teacher's.

Post-hoc explainability methods such as \textbf{LIME} \cite{ribeiro2016should} and \textbf{SHAP} \cite{lundberg2017unified} treat the model as a black box, while gradient-based methods such as \textbf{Integrated Gradients} \cite{sundararajan2017axiomatic} exploit differentiability; all three are commonly used to quantify token-level contributions to a prediction. Beyond post-hoc attribution, attention maps offer a potentially interpretable internal mechanism, but whether they constitute a faithful explanation remains contested: Jain and Wallace \cite{jain2019attention} show that alternative attention configurations can induce similar predictions, whereas Wiegreffe and Pinter \cite{wiegreffe2019attention} argue attention can still provide plausible explanations under appropriate conditions.

\section{Experiments}

\subsection{Experimental Setup}\label{sec:setup}

The project used \textbf{Python}, \textbf{PyTorch}, Hugging Face's \textbf{transformers}, \textbf{lime}/\textbf{shap} for explanation, \textbf{SciPy} for the rank/value-correlation statistics used in the aggregate explanation metrics, and \textbf{scikit-learn} for evaluation metrics.

The teacher is \textbf{TinyLlama-1.1B-Chat-v1.0} \cite{zhang2024tinyllama}, a $22$-layer decoder-only Transformer \cite{vaswani2017attention}. Every student shares the teacher's embeddings, attention mechanism, and hidden dimensions, reducing only depth to $11$ randomly-initialized blocks ($2{:}1$ Transformer layer correspondence; compression rate $CR\approx1.787$, $\approx44\%$).

The studied problem is reformulated as a generative classification \cite{schick2021exploiting}, where class-labels are single vocabulary tokens. Dataset inputs are inserted into a fixed prompt template, and only the class-token logits are compared at inference. For \textbf{IMDb} \cite{maas2011learning} ($50000$ reviews, labels \textit{positive}/\textit{negative}), the template is \texttt{"Review: \{input\} Summary:"}; filtered to the teacher's $2048$-token context, splits are $22492$/$2500$/$24995$ (train/val/test). For \textbf{AG News} \cite{zhang2015character} ($4$ classes: \textit{World}, \textit{Sports}, \textit{Business}, \textit{Sci/Tech}), the template is \texttt{"Article: \{input\} Theme:"}. Since \textit{Sci/Tech} tokenizes to four sub-tokens, it is relabeled \textit{Science} (one token) purely to preserve the single-token framing. Splits are $30000$/$100$/$7600$ (balanced). All other hyperparameters are shared across datasets.

The teacher is fully fine-tuned per dataset for one epoch (learning rate $2\times10^{-5}$, batch size $1$, AdamW, early stopping after $5$ non-improving validation rounds). Every student is randomly initialized and trained for $3$ epochs (learning rate $1\times10^{-4}$, no early stopping), minimizing \eqref{eq:loss}.

\begin{align}
\mathcal{L} &= \alpha\,\mathcal{L}_{CE} + (1-\alpha)\,T^2 D_{KL}\left(P_P \Vert P_S\right) \notag \\
&\quad + \beta\,T^2 D_{KL}\left(P_P^{(c)} \Vert P_S^{(c)}\right)
\label{eq:loss}
\end{align}

with $T=2$ fixed, $P_P^{(c)},P_S^{(c)}$ the probability distributions restricted to the class-label tokens, and $(\alpha,\beta)$ set per Table~\ref{tab:model_family}, defining a family of four students.

\begin{table}[bp]
\centering
\caption{Model family. All students share the same reduced ($11$-layer) architecture and differ only in $(\alpha,\beta)$.}
\small
\begin{tabular}{|l|c|c|c|l|}
\hline
Model & Init. & $\alpha$ & $\beta$ & Description \\ \hline
$P$ (teacher) & pretrained & -- & -- & fine-tuned, CE only \\
$B$ & random & $1$ & $0$ & no distillation \\
$S_0$ & random & $0.9$ & $0$ & weak global $KD$ \\
$S_1$ & random & $0.5$ & $0$ & standard global $KD$ \\
$S_2$ & random & $0.5$ & $0.25$ & $+$ class-restricted $KD$ \\ \hline
\end{tabular}
\label{tab:model_family}
\end{table}

Explanations from \textbf{LIME}, \textbf{SHAP}, and \textbf{Integrated Gradients} (zero-embedding baseline) are compared between teacher and student via aggregate metrics, computed over $N=100$ held-out examples per dataset for all three methods. Representational comparisons use the $2{:}1$ teacher-student layer correspondence: attention maps are averaged over held-out samples ($N=200$) at each layer, and student layer $i$ is then compared against the average of teacher layers $2i$ and $2i+1$. All experiments ran on an NVIDIA RTX 5070 Ti (16GB).

\subsection{Comparison Metrics}\label{sec:metrics}

We compare Attention maps with complementary geometric, probabilistic, and spectral measures. Given two attention matrices $A,B$ and their vectorizations $\vec A, \vec B$, we compared them via \textit{cosine similarity} $\cos(\theta)=\langle \vec A,\vec B\rangle/(\|\vec A\|_2\|\vec B\|_2)\in[-1,1]$, the \textit{Frobenius distance} $d_F(A,B)=\|A-B\|_F=\sum_{i,j}(A-B)^2_{ij}$ and the row-averaged \textit{Jensen--Shannon divergence} $\frac{1}{n}\sum_i JS(A_i||B_i)$ were $A_i$ is the $i_{th}$ row of $A$. Since Attention maps are $n \times n$ row-stochastic matrices, the \textit{Frobenius distance} is bounded by $\sqrt{2n}$ ($64$ for our $n=2048$ maps). The \textit{Jensen--Shannon divergence} is theoretically bounded in $[0,\ln 2]$ if measured in nats \cite{lin1991divergence}, where $0$ means equal distributions. For spectral metrics we propose the entropy-based \textit{effective rank} \cite{roy2007effective} and the \textit{Participation Ratio} \cite{recanatesi2022scale}.

For evaluating explainability methods, we compute aggregate metrics over held-out samples, comparing each student's word-level attribution against the teacher's for every method. Metrics include top-$k$ \textit{Jaccard overlap} ($k=10$), \textit{Spearman correlation} of signed scores over the common attributed vocabulary, and deletion-based \textit{faithfulness@$k$} \cite{deyoung2020eraser}, the drop in true-class probability after masking the top-$k$ words (computed per model, including the teacher, as reference). Given \textbf{LIME's} stochastic nature, we also measure \textit{stability} as the Jaccard@10 mean across $5$ re-runs with different seeds.

\subsection{Evaluation on the Test Set}\label{sec:test-eval}

Table~\ref{tab:perf} reports classification performance and average $KL$ divergence to the teacher ($T=1,2$) for all models on both datasets. On \textbf{IMDb}, distillation's effect on accuracy is mixed: only the class-restricted $S_2$ edges out the non-distilled $B$ ($0.867$ vs.\ $0.858$), while $S_0$ and $S_1$ both \emph{underperform} $B$ ($0.813$ and $0.798$), i.e.\ the distillation term hurts relative to training from scratch. On \textbf{AG News}, this pattern is more pronounced: $B$ achieves the \emph{highest} accuracy of all five models ($B>S_0>S_2>S_1>P$). $KL$ fidelity, by contrast, follows the same monotonic pattern on \emph{both} datasets and temperatures: $B$ is by far the least faithful, improving with distillation strength ($B>S_0>S_1\gtrsim S_2$). Notably, on \textbf{IMDb} $S_1$ attains \emph{lower} $KL$ than $S_2$ despite worse accuracy: global $KL$ divergence is not a predictive proxy for classification performance, nor, as shown below, of explanatory or representational alignment.

\begin{table}[tb]
\centering
\caption{Accuracy / F1 (macro-averaged for \textbf{AG News}) on the test set, and average $KL$ divergence to the teacher $P$ at $T=1,2$ (\textbf{IMDb}: $N=3000$ balanced subset; \textbf{AG News}: full test set).}
\small
\begin{tabular}{|l|l|c|c|c|c|c|}
\hline
Dataset & Model & Acc. & F1 & $KL_{T=1}$ & $KL_{T=2}$ \\ \hline
\textbf{IMDb} & $P$ & $0.95$ & $0.95$ & -- & -- \\
\textbf{IMDb} & $B$ & $0.86$ & $0.86$ & $2.85$ & $1.69$ \\
\textbf{IMDb} & $S_0$ & $0.81$ & $0.81$ & $0.10$ & $0.13$ \\
\textbf{IMDb} & $S_1$ & $0.80$ & $0.80$ & $0.10$ & $0.13$ \\
\textbf{IMDb} & $S_2$ & $0.87$ & $0.87$ & $0.11$ & $0.19$ \\ \hline
\textbf{AG News} & $P$ & $0.85$ & $0.85$ & -- & -- \\
\textbf{AG News} & $B$ & $0.89$ & $0.89$ & $2.08$ & $1.78$ \\
\textbf{AG News} & $S_0$ & $0.88$ & $0.88$ & $0.63$ & $0.42$ \\
\textbf{AG News} & $S_1$ & $0.88$ & $0.88$ & $0.28$ & $0.23$ \\
\textbf{AG News} & $S_2$ & $0.88$ & $0.88$ & $0.25$ & $0.25$ \\ \hline
\end{tabular}
\label{tab:perf}
\end{table}

\subsection{Functional Evaluation and Explanation Comparison}\label{sec:functional-eval}

\textbf{LIME} explanations for a representative \textbf{AG News} \textit{Business} example show the teacher's top attributions concentrated on topical terms (\textit{pension}, \textit{workers}, \textit{Unions}, \textit{Federal}); students share most of these words, mixing in generic phrase tokens absent from the teacher's list. A similar pattern holds on \textbf{IMDb}, where the teacher and most students concentrate on genuine sentiment words (\textit{extraordinary}, \textit{wonderful}, \textit{dull}, \textit{clichéd}), except $S_1$, which consistently favors function words (\textit{it}, \textit{and}, \textit{the}, \textit{to}) over content words regardless of its classification performance.

\textbf{SHAP} explanations for the same examples are far less concentrated: the top-$10$ words capture only a small fraction of the total attribution, with the remaining $\sim$140 tokens accounting for most of it, and content/function words appear interspersed across all five models rather than following a clean teacher/student split. \textbf{Integrated Gradients}, for every model including the teacher, is instead dominated by one or two extreme-magnitude spikes on semantically arbitrary tokens (proper nouns, subword fragments) rather than topical or sentiment content.

Table~\ref{tab:stability} reports \textbf{LIME} stability (mean pairwise Jaccard@10 across $5$ re-runs, $N=20$). The teacher is more stable than the students on \textbf{AG News} but not always on \textbf{IMDb}, so stability alone is not a proxy for explanation quality.

Table~\ref{tab:explanation_metrics} reports aggregate \textbf{LIME} agreement (vs.\ teacher) and faithfulness on both datasets, $N=100$, $k=10$. On both, $S_2$ is consistently the best-aligned and most-faithful student, though the gap to $B$ is modest.

\textbf{SHAP} (Table~\ref{tab:explanation_metrics}) shows the same $S_2$-best ranking, but faithfulness near zero for \emph{every} model, including the teacher -- likely a limitation of masking-based perturbation rather than a genuine absence of faithful explanations.

\textbf{Integrated Gradients} (Table~\ref{tab:explanation_metrics}) shows high top-$10$ token overlap ($\approx0.78$--$0.89$) alongside near-zero or negative rank correlations, indicating agreement on \emph{which} tokens matter without agreement on their relative importance. Sign agreement@$10$ (not tabulated) mirrors this split: \textbf{LIME}/\textbf{SHAP} follow the same $S_2$-best pattern ($0.72$--$0.82$), whereas IG sits at chance level ($0.50$--$0.54$) for every model, so its cross-model consistency is comparatively shallow.

\begin{table}[bp]
\centering
\caption{\textbf{LIME} stability (mean pairwise Jaccard@10 across $5$ seeds, $N=20$), both datasets.}
\small
\begin{tabular}{|l|c|c|c|c|c|}
\hline
Dataset & $P$ & $B$ & $S_0$ & $S_1$ & $S_2$ \\ \hline
\textbf{IMDb} & $0.55$ & $0.59$ & $0.68$ & $0.61$ & $0.71$ \\
\textbf{AG News} & $0.77$ & $0.67$ & $0.64$ & $0.62$ & $0.69$ \\ \hline
\end{tabular}
\label{tab:stability}
\end{table}

\subsection{Representational Comparison}\label{sec:repr-comparison}

Table~\ref{tab:repr_metrics} reports cosine similarity, Frobenius distance, and Jensen--Shannon divergence between each student's attention map and the teacher's average ($P_{avg}$, matched via the $2{:}1$ correspondence), across six layers per dataset (full $11$-layer tables are in the extended version). On \textbf{IMDb}, $B$ is slightly lower than the three distilled variants at most layers by cosine similarity, some evidence for attention alignment with the teacher. In \textbf{AG News}, however, the pattern is reversed: $B$ is comparable to, and at several layers more similar to, the teacher than any distilled variant. Both datasets thus suggest attention-level similarity to the teacher is not a reliable signature of distillation, since the non-distilled baseline is nearly as close, or closer, to the teacher's attention geometry.

\begin{table*}[tb]
\centering
\caption{\textbf{LIME}, \textbf{SHAP}, and \textbf{Integrated Gradients} (IG, zero baseline) Jaccard@10, Spearman $\rho$, and Faithfulness@10 vs.\ teacher, both datasets ($N=100$, $k=10$).}
\footnotesize
\begin{tabular}{|l|l|c|c|c|c|c|c|c|c|c|}
\hline
\multirow{2}{*}{Dataset} & \multirow{2}{*}{Model} & \multicolumn{3}{c|}{LIME} & \multicolumn{3}{c|}{SHAP} & \multicolumn{3}{c|}{IG} \\
\cline{3-11}
 & & Jac.@10 & $\rho$ & Faithf. & Jac.@10 & $\rho$ & Faithf. & Jac.@10 & $\rho$ & Faithf. \\ \hline
\textbf{IMDb} & $P$ & -- & -- & $0.17$ & -- & -- & $0.12$ & -- & -- & $0.01$ \\
\textbf{IMDb} & $B$ & $0.18$ & $0.25$ & $0.08$ & $0.12$ & $0.24$ & $0.03$ & $0.79$ & $-0.01$ & $0.01$ \\
\textbf{IMDb} & $S_0$ & $0.17$ & $0.22$ & $0.20$ & $0.12$ & $0.24$ & $0.09$ & $0.78$ & $0.03$ & $0.09$ \\
\textbf{IMDb} & $S_1$ & $0.14$ & $0.19$ & $0.18$ & $0.10$ & $0.25$ & $0.06$ & $0.81$ & $0.01$ & $0.03$ \\
\textbf{IMDb} & $S_2$ & $0.21$ & $0.30$ & $0.27$ & $0.18$ & $0.33$ & $0.12$ & $0.79$ & $0.04$ & $0.02$ \\
\hline
\textbf{AG News} & $P$ & -- & -- & $0.14$ & -- & -- & $0.00$ & -- & -- & $0.05$ \\
\textbf{AG News} & $B$ & $0.34$ & $0.32$ & $0.31$ & $0.32$ & $0.27$ & $-0.00$ & $0.89$ & $0.02$ & $0.07$ \\
\textbf{AG News} & $S_0$ & $0.35$ & $0.28$ & $0.40$ & $0.32$ & $0.28$ & $0.00$ & $0.88$ & $0.04$ & $0.16$ \\
\textbf{AG News} & $S_1$ & $0.34$ & $0.31$ & $0.30$ & $0.32$ & $0.27$ & $0.00$ & $0.88$ & $-0.04$ & $0.14$ \\
\textbf{AG News} & $S_2$ & $0.34$ & $0.34$ & $0.43$ & $0.31$ & $0.26$ & $-0.00$ & $0.89$ & $0.01$ & $0.34$ \\ \hline
\end{tabular}
\label{tab:explanation_metrics}
\end{table*}

\begin{table*}[tb]
\centering
\caption{Cosine similarity, Frobenius distance, and Jensen--Shannon divergence between each student's attention map and the teacher's $P_{avg}$, across layers $i$ (Frobenius theoretical maximum $\sqrt{2n}=64$ for $n=2048$).}
\footnotesize
\begin{tabular}{|l|c|c|c|c|c|c|c|c|c|c|c|c|c|}
\hline
\multirow{2}{*}{Dataset} & \multirow{2}{*}{$i$} & \multicolumn{4}{c|}{Cosine} & \multicolumn{4}{c|}{Frobenius} & \multicolumn{4}{c|}{$JS$} \\
\cline{3-14}
 & & $B$ & $S_0$ & $S_1$ & $S_2$ & $B$ & $S_0$ & $S_1$ & $S_2$ & $B$ & $S_0$ & $S_1$ & $S_2$ \\ \hline
\textbf{IMDb} & $0$ & $0.93$ & $0.92$ & $0.94$ & $0.92$ & $1.27$ & $1.36$ & $1.24$ & $1.35$ & $0.01$ & $0.018$ & $0.011$ & $0.02$ \\
\textbf{IMDb} & $2$ & $0.29$ & $0.42$ & $0.46$ & $0.63$ & $36.28$ & $35.66$ & $35.44$ & $33.95$ & $0.43$ & $0.39$ & $0.38$ & $0.34$ \\
\textbf{IMDb} & $3$ & $0.23$ & $0.43$ & $0.50$ & $0.69$ & $39.12$ & $38.21$ & $37.84$ & $36.46$ & $0.49$ & $0.44$ & $0.43$ & $0.39$ \\
\textbf{IMDb} & $4$ & $0.15$ & $0.37$ & $0.42$ & $0.65$ & $26.46$ & $25.50$ & $25.08$ & $22.95$ & $0.38$ & $0.31$ & $0.30$ & $0.26$ \\
\textbf{IMDb} & $7$ & $0.13$ & $0.26$ & $0.48$ & $0.38$ & $18.04$ & $17.39$ & $16.08$ & $16.81$ & $0.37$ & $0.27$ & $0.24$ & $0.25$ \\
\textbf{IMDb} & $10$ & $0.14$ & $0.27$ & $0.23$ & $0.16$ & $21.31$ & $20.66$ & $20.85$ & $21.18$ & $0.41$ & $0.30$ & $0.31$ & $0.32$ \\
\hline
\textbf{AG News} & $0$ & $0.98$ & $0.97$ & $0.97$ & $0.98$ & $1.08$ & $1.46$ & $1.39$ & $1.22$ & $0.01$ & $0.01$ & $0.01$ & $0.01$ \\
\textbf{AG News} & $2$ & $0.43$ & $0.32$ & $0.23$ & $0.37$ & $39.22$ & $40.16$ & $40.64$ & $39.64$ & $0.45$ & $0.48$ & $0.50$ & $0.47$ \\
\textbf{AG News} & $3$ & $0.44$ & $0.39$ & $0.22$ & $0.23$ & $37.78$ & $38.51$ & $39.48$ & $39.45$ & $0.42$ & $0.43$ & $0.47$ & $0.47$ \\
\textbf{AG News} & $4$ & $0.52$ & $0.35$ & $0.29$ & $0.35$ & $21.42$ & $22.96$ & $23.26$ & $22.91$ & $0.25$ & $0.23$ & $0.23$ & $0.22$ \\
\textbf{AG News} & $7$ & $0.27$ & $0.34$ & $0.33$ & $0.28$ & $18.05$ & $17.52$ & $17.57$ & $17.91$ & $0.31$ & $0.23$ & $0.22$ & $0.25$ \\
\textbf{AG News} & $10$ & $0.31$ & $0.27$ & $0.23$ & $0.26$ & $19.33$ & $19.52$ & $19.82$ & $19.62$ & $0.30$ & $0.24$ & $0.26$ & $0.25$ \\ \hline
\end{tabular}
\label{tab:repr_metrics}
\end{table*}

By contrast, Frobenius distance and $JS$ divergence (Table~\ref{tab:repr_metrics}) are nearly indistinguishable across all four models at every layer on both datasets, following the same profile: near-perfect alignment at $i=0$, peaking at $i=2$--$3$ (Frobenius $34$--$41$ out of a theoretical maximum of $64$; $JS$ $0.34$--$0.50$), then stabilizing by $i=4$--$10$. At an even finer grain, per-head attention-similarity analysis shows $B$ numerically indistinguishable from the three distilled students on \emph{both} datasets -- the clearest evidence that per-head similarity to the teacher is insensitive to whether a model was distilled.

Table~\ref{tab:spectral} reports Participation Ratio ($PR$) and entropy-based effective rank ($ER$) on both datasets, at the same layers as Table~\ref{tab:repr_metrics}. On both measures, $P_{avg}$ stays close to rank-one throughout, while the four students show no consistent ranking: on \textbf{AG News}, $B$ trends from lowest to highest as depth increases on both $PR$ and $ER$; on \textbf{IMDb}, the relative order among $B$, $S_0$, $S_1$, and $S_2$ instead shifts from layer to layer. Spectral complexity is therefore not obviously ordered by distillation strength on either dataset.

\subsection{Discussion}\label{sec:discussion}

Taken together, these results indicate that, at least in the studied settings, standard KD does not directly induce alignment of either explanatory mechanisms or internal representations between teacher and student. Both properties instead appear to emerge from the shared dataset and task rather than from the distillation term itself. Distillation's only clearly attributable effect is distributional: it reliably pulls the student's output distribution toward the teacher's (Table~\ref{tab:perf}). Accuracy itself is barely affected by distillation, and inconsistently so across datasets -- but these remain well-behaved, competitive classifiers, worth studying on those terms alone.

This dissociation is not incidental: if explanatory or representational alignment were induced by the distillation term itself, both should track distillation strength (i.e.\ $KL$ fidelity) directly. On \textbf{IMDb}, $S_1$ attains \emph{lower} $KL$ than $S_2$ (Table~\ref{tab:perf}) yet is markedly less \textbf{LIME}/\textbf{SHAP}-aligned with the teacher, the opposite of what a distillation-driven account would predict; \textbf{Integrated Gradients} shows no consistent best-aligned student at all, despite the same spread in $KL$ fidelity. On \textbf{AG News}, $B$ shows no disadvantage, and at several layers an advantage, in attention-level similarity to the teacher (Table~\ref{tab:repr_metrics}), despite having the highest accuracy of all five models. Neither similarity scales with distillation strength, reinforcing that the shared task and data are the more likely source of teacher-student similarity.

\subsection{Does Distillation Add Value Beyond Depth Reduction?}\label{sec:distillation-value}

Distillation's only robust, task-independent effect is pulling the student's output distribution toward the teacher's. Its effects on accuracy, explanatory alignment, and fine-grained representational similarity are comparatively small, inconsistent across datasets, and, for attention geometry, largely absent once a non-distilled baseline is available for comparison.

\begin{table*}[tb]
\centering
\caption{Participation Ratio ($PR$) and entropy-based effective rank ($ER$), at layers $i$, both datasets.}
\footnotesize
\begin{tabular}{|l|c|c|c|c|c|c|c|c|c|c|c|}
\hline
\multirow{2}{*}{Dataset} & \multirow{2}{*}{$i$} & \multicolumn{5}{c|}{$PR$} & \multicolumn{5}{c|}{$ER$} \\
\cline{3-12}
 & & $B$ & $S_0$ & $S_1$ & $S_2$ & $P_{avg}$ & $B$ & $S_0$ & $S_1$ & $S_2$ & $P_{avg}$ \\ \hline
\textbf{IMDb} & $0$ & $2.6$ & $2.6$ & $2.5$ & $2.7$ & $3.3$ & $60$ & $117$ & $96$ & $107$ & $309$ \\
\textbf{IMDb} & $2$ & $3.4$ & $3.0$ & $2.5$ & $1.7$ & $1.0$ & $278$ & $298$ & $230$ & $180$ & $5$ \\
\textbf{IMDb} & $3$ & $2.9$ & $3.3$ & $2.5$ & $2.0$ & $1.0$ & $328$ & $398$ & $304$ & $361$ & $7$ \\
\textbf{IMDb} & $4$ & $5.0$ & $4.4$ & $4.3$ & $2.4$ & $1.0$ & $315$ & $437$ & $517$ & $485$ & $20$ \\
\textbf{IMDb} & $7$ & $13.7$ & $13.4$ & $7.5$ & $7.1$ & $1.0$ & $505$ & $555$ & $505$ & $488$ & $84$ \\
\textbf{IMDb} & $10$ & $9.7$ & $8.6$ & $17.0$ & $12.5$ & $1.0$ & $491$ & $509$ & $551$ & $528$ & $129$ \\
\hline
\textbf{AG News} & $0$ & $1.3$ & $1.3$ & $1.4$ & $1.4$ & $1.4$ & $13$ & $14$ & $16$ & $18$ & $24$ \\
\textbf{AG News} & $2$ & $1.6$ & $1.4$ & $1.7$ & $2.0$ & $1.0$ & $29$ & $14$ & $31$ & $39$ & $1$ \\
\textbf{AG News} & $3$ & $1.9$ & $1.4$ & $1.9$ & $2.1$ & $1.0$ & $40$ & $19$ & $32$ & $51$ & $2$ \\
\textbf{AG News} & $4$ & $1.9$ & $1.4$ & $1.8$ & $1.9$ & $1.0$ & $42$ & $20$ & $47$ & $55$ & $3$ \\
\textbf{AG News} & $7$ & $3.4$ & $1.5$ & $1.8$ & $1.6$ & $1.0$ & $62$ & $32$ & $47$ & $47$ & $5$ \\
\textbf{AG News} & $10$ & $3.3$ & $1.6$ & $1.7$ & $1.8$ & $1.0$ & $66$ & $28$ & $40$ & $43$ & $6$ \\ \hline
\end{tabular}
\label{tab:spectral}
\end{table*}

\subsection{Limitations}\label{sec:limitations}

This study covers only two English text-classification datasets of modest label cardinality; whether the task-dependence between them generalizes to other modalities or higher-cardinality tasks is untested. The $(\alpha,\beta)$ grid is narrow: only two points beyond $S_1$/$S_2$ were tested ($\alpha=1$ for $B$, $\alpha=0.9$ for $S_0$), and the opposite extreme (pure distillation, $\alpha=0$) was never tested. Underlying all of this is a hardware constraint: every experiment ran on a single consumer GPU (RTX 5070 Ti, $16$GB), bounding how many datasets, scales, and architectures could be explored.

\section{Conclusions}

This work investigated whether functional similarity induced by knowledge distillation extends to internal representations and explanatory mechanisms, using a family of models independently on two tasks: binary \textbf{IMDb} sentiment classification and four-class \textbf{AG News} topic classification. In these comparatively simple settings, the answer is negative rather than definitively so: every student, including the baseline, assigns importance to irrelevant tokens regardless of accuracy, and $B$ is indistinguishable from the distilled students at the level of individual attention heads on both datasets; even at the level of full attention maps, distillation offers at best weak, inconsistent evidence of an effect on \textbf{IMDb} and none on \textbf{AG News}. The one property that behaves consistently, independently of accuracy, is distributional ($KL$) fidelity to the teacher.

A striking pattern underlies this: $B$, $S_0$, $S_1$, and $S_2$ converge toward similar attention structure despite spanning $\alpha$ from $0.5$ to $1$. All four share the same architecture, dataset, and cross-entropy training, with only $S_0$, $S_1$, and $S_2$ receiving a distillation term. Since $B$ still converges to the same structure as the distilled students, that convergence cannot be attributed to distillation alone. A natural next step, left for future work, is a student trained with \emph{pure} distillation ($\alpha=0$, no cross-entropy term): if its representations diverge from the $B$/$S_0$/$S_1$/$S_2$ cluster, that would implicate cross-entropy over the dataset, rather than distillation, as the source of this convergence.

\section*{Acknowledgment}

The writing of this manuscript was partially assisted by generative language models used exclusively for stylistic revision and textual reformulation tasks: Claude Sonnet 4.6 \cite{anthropic2025claude} (Anthropic).

\bibliographystyle{plain}
\bibliography{references}

\end{document}
